\documentclass[aspectratio=169,9pt]{beamer}

\usetheme{default}
\usefonttheme{professionalfonts}
\setbeamertemplate{navigation symbols}{}
\setbeamertemplate{footline}{\hfill\scriptsize\insertframenumber/3\hspace{2mm}\vspace{1mm}}

\usepackage{amsmath,amssymb,bm,mathtools}
\usepackage{xcolor}
\usepackage[most]{tcolorbox}

\definecolor{deepblue}{RGB}{0,70,150}
\definecolor{softblue}{RGB}{232,242,255}
\definecolor{softgreen}{RGB}{235,250,238}
\definecolor{softorange}{RGB}{255,245,230}
\definecolor{softpurple}{RGB}{245,238,255}
\definecolor{darkred}{RGB}{160,35,25}
\definecolor{darkgreen}{RGB}{20,120,60}
\definecolor{darkpurple}{RGB}{95,45,150}

\setbeamerfont{frametitle}{series=\bfseries,size=\Large}
\setbeamercolor{frametitle}{fg=black}
\setbeamertemplate{itemize item}{\color{deepblue}$\blacktriangleright$}
\setlength{\abovedisplayskip}{3pt}
\setlength{\belowdisplayskip}{3pt}
\setlength{\jot}{1pt}

\newtcolorbox{bluebox}[1][]{colback=softblue,colframe=deepblue,boxrule=0.8pt,arc=1.5mm,left=1.5mm,right=1.5mm,top=0.8mm,bottom=0.8mm,#1}
\newtcolorbox{greenbox}[1][]{colback=softgreen,colframe=darkgreen,boxrule=0.8pt,arc=1.5mm,left=1.5mm,right=1.5mm,top=0.8mm,bottom=0.8mm,#1}
\newtcolorbox{orangebox}[1][]{colback=softorange,colframe=darkred,boxrule=0.8pt,arc=1.5mm,left=1.5mm,right=1.5mm,top=0.8mm,bottom=0.8mm,#1}
\newtcolorbox{purplebox}[1][]{colback=softpurple,colframe=darkpurple,boxrule=0.8pt,arc=1.5mm,left=1.5mm,right=1.5mm,top=0.8mm,bottom=0.8mm,#1}

\newcommand{\Evec}{\bm{E}}
\newcommand{\Jvec}{\bm{J}}
\newcommand{\Jn}{\bm{J}_n}
\newcommand{\Jp}{\bm{J}_p}
\newcommand{\qrad}{\dot{q}_{\mathrm{rad}}}
\newcommand{\qth}{\dot{q}_{\mathrm{th}}}
\newcommand{\dd}{\,\mathrm{d}}
\newcommand{\Rsrh}{R_{\mathrm{SRH}}}

\begin{document}

\begin{frame}{Module 5: Carrier conservation and drift--diffusion transport}
\footnotesize
\vspace{-4mm}
\begin{bluebox}
\centering\small
\textbf{Goal:} convert electrostatic fields, defect maps, and temperature maps into mobile-carrier densities, current densities, and measurable electrical degradation.
\end{bluebox}

\vspace{1mm}
\begin{columns}[T,totalwidth=\textwidth]
\begin{column}{0.50\textwidth}
\textbf{Carrier conservation}
\begin{equation*}
\boxed{\frac{\partial n}{\partial t}=\frac{1}{q}\nabla\cdot\Jn+G_n-R_n}
\end{equation*}
\begin{equation*}
\boxed{\frac{\partial p}{\partial t}=-\frac{1}{q}\nabla\cdot\Jp+G_p-R_p}
\end{equation*}
\small
Here \(n\) and \(p\) are electron and hole densities, \(q\) is the elementary charge, \(G\) is generation, and \(R\) is recombination.

\vspace{1mm}
The electric field is supplied by Module 2:
\begin{equation*}
\boxed{\Evec=-\nabla \phi}
\end{equation*}
\end{column}

\begin{column}{0.48\textwidth}
\textbf{Drift--diffusion current densities}
\begin{equation*}
\boxed{\Jn=q\mu_n n\Evec+qD_n\nabla n}
\end{equation*}
\begin{equation*}
\boxed{\Jp=q\mu_p p\Evec-qD_p\nabla p}
\end{equation*}
\small
Electrons drift opposite \(\Evec\), but electron \emph{conventional current} points along \(\Evec\). Holes drift and carry conventional current along \(\Evec\).

\vspace{1mm}
Near nondegenerate equilibrium,
\begin{equation*}
\boxed{D_n=\frac{k_B T}{q}\mu_n},\qquad
\boxed{D_p=\frac{k_B T}{q}\mu_p}
\end{equation*}
\small
where \(D_n,D_p\) are diffusivities and \(\mu_n,\mu_p\) are mobilities.
\end{column}
\end{columns}

\vspace{1mm}
\begin{orangebox}
\centering\scriptsize
\textbf{Main physics:} field gradients steer carriers; concentration gradients spread carriers; defects reduce mobility and remove carriers through trap-assisted recombination.
\end{orangebox}
\end{frame}

\begin{frame}{Module 5: Defects as mobility modifiers and recombination centers}
\footnotesize
\vspace{-4mm}
\begin{bluebox}
\centering\small
Radiation-induced defects enter Module 5 in two distinct ways: they scatter mobile carriers and they introduce trap states that enable electron-hole recombination.
\end{bluebox}

\vspace{1mm}
\begin{columns}[T,totalwidth=\textwidth]
\begin{column}{0.50\textwidth}
\textbf{Defect-limited mobility}
\begin{equation*}
\boxed{\mu_n(T,\{C_j\})=
\frac{\mu_{n0}(T)}{1+\sum_j \alpha_{n,j}C_j}}
\end{equation*}
\begin{equation*}
\boxed{\mu_p(T,\{C_j\})=
\frac{\mu_{p0}(T)}{1+\sum_j \alpha_{p,j}C_j}}
\end{equation*}
\small
\(C_j\) is the concentration of defect species \(j\) from Module 3.  The parameters \(\alpha_{n,j}\) and \(\alpha_{p,j}\) summarize how strongly that defect scatters electrons and holes.

\vspace{1mm}
\begin{equation*}
\boxed{\mu_{n0}(T)=\mu_{n,\mathrm{ref}}\left(\frac{T}{T_{\mathrm{ref}}}\right)^{-\gamma_n}}
\end{equation*}
\small
and similarly for holes.  This keeps the reduced model monotonic: more scattering defects \(\Rightarrow\) lower mobility.
\end{column}

\begin{column}{0.48\textwidth}
\textbf{Shockley--Read--Hall recombination}
\begin{equation*}
\boxed{\Rsrh=\frac{np-n_i^2}
{\tau_p(n+n_1)+\tau_n(p+p_1)}}
\end{equation*}
\begin{equation*}
\boxed{\tau_n=\frac{1}{\sigma_n v_{\mathrm{th}}N_t}},\qquad
\boxed{\tau_p=\frac{1}{\sigma_p v_{\mathrm{th}}N_t}}
\end{equation*}
\small
\(N_t\) is the active trap density, \(\sigma_n,\sigma_p\) are capture cross sections, and \(v_{\mathrm{th}}\) is thermal velocity.

\vspace{1mm}
For multiple radiation-induced trap species,
\begin{equation*}
\boxed{R_{\mathrm{def}}(x,y,t)=\sum_j
R_{\mathrm{SRH},j}\!ig(n,p,T,C_j\big)}
\end{equation*}
\small
Each electrically active defect species can contribute a recombination pathway.
\end{column}
\end{columns}

\vspace{1mm}
\begin{purplebox}
\centering\scriptsize
\textbf{Interpretation:} the same defect map can lower current by reducing \(\mu_n,\mu_p\), increasing \(R_{\mathrm{def}}\), and changing space charge through Module 2.
\end{purplebox}
\end{frame}


\begin{frame}{Module 5: Coupled reduced PDE and device response}
\scriptsize
\vspace{-3mm}
\begin{bluebox}
\centering
\textbf{Practical role:} solve carrier transport using fields from Modules 1--4, then reduce the solution to current, leakage, and device-response metrics.
\end{bluebox}

\vspace{0.8mm}
\begin{columns}[T,totalwidth=\textwidth]
\begin{column}{0.52\textwidth}
\textbf{Coupled 2D carrier equations}
\begin{equation*}
\boxed{
\frac{\partial n}{\partial t}
=\frac{1}{q}\nabla\cdot\left(q\mu_n n\Evec+qD_n\nabla n\right)+G_n-R_{\mathrm{def}}
}
\end{equation*}
\begin{equation*}
\boxed{
\frac{\partial p}{\partial t}
=-\frac{1}{q}\nabla\cdot\left(q\mu_p p\Evec-qD_p\nabla p\right)+G_p-R_{\mathrm{def}}
}
\end{equation*}

\vspace{0.4mm}
\textbf{Coefficient dependence}
\begin{equation*}
\boxed{\mu_{n,p}=\mu_{n,p}(T,C_1,C_2,\ldots)},
\qquad
\boxed{R_{\mathrm{def}}=R_{\mathrm{def}}(n,p,T,C_j)}
\end{equation*}
\smallskip
\begin{orangebox}
\textbf{Interpretation:} defects change the PDE coefficients and the sink terms; temperature changes mobility, diffusivity, and trap kinetics.
\end{orangebox}
\end{column}

\begin{column}{0.46\textwidth}
\textbf{Module handoffs and outputs}
\begin{equation*}
\boxed{G_{n,p}^{\mathrm{rad}}\simeq
\eta_{eh}\frac{\qrad^{2D}}{\varepsilon_{eh}}}
\end{equation*}
\small
Radiation energy deposition can generate electron-hole pairs through the fraction \(\eta_{eh}\).

\vspace{0.5mm}
\begin{equation*}
\boxed{\rho=q\left(p-n+N_D^+-N_A^-+\sum_j z_j C_j\right)}
\end{equation*}
\begin{equation*}
\boxed{\nabla\cdot(\epsilon\nabla\phi)=-\rho},
\qquad
\boxed{\Evec=-\nabla\phi}
\end{equation*}

\vspace{0.5mm}
\begin{equation*}
\boxed{I_c(t)=\int_{\Gamma_c}(\Jn+\Jp)\cdot\hat{\bm{n}}\,\dd s}
\end{equation*}
\end{column}
\end{columns}

\vspace{0.8mm}
\begin{greenbox}
\centering\scriptsize
\textbf{Module 5 output:} \(n,p,\Jn,\Jp\), recombination loss, contact currents, leakage/current shifts, and carrier maps for Module 6.
\end{greenbox}
\end{frame}

\end{document}
